% TOP_PROBLEM: {"id": 30006801, "title": "Berry random-wave conjecture: real negative-curvature BS formulation", "category_id": 16, "category": {"id": 16, "name": "physics", "display_name": "Mathematical Physics", "description": "Problems at the intersection of mathematics and physics.", "slug": "physics", "order_index": 16, "created_at": "2026-07-31T15:26:25.671Z"}, "status": "open"}
% FIRST_POSTED: 2026-09-27
% !TeX program = lualatex
% Compile twice: lualatex 30006801.tex
% All references and prose are embedded; no BibTeX database or external inputs.
\documentclass[11pt,a4paper]{article}
\usepackage[margin=24mm,headheight=14pt]{geometry}
\usepackage{fontspec}
\setmainfont{Latin Modern Roman}
\setsansfont{Latin Modern Sans}
\setmonofont{Latin Modern Mono}
\usepackage{amsmath,amssymb,mathtools}
\usepackage{unicode-math}
\setmathfont{Latin Modern Math}
\usepackage{microtype}
\usepackage{enumitem,xurl,needspace}
\usepackage[dvipsnames]{xcolor}
\usepackage{hyperref}
\hypersetup{unicode=true,colorlinks=true,linkcolor=MidnightBlue,urlcolor=MidnightBlue,
 pdftitle={Research Notebook: Section 30006801},pdfauthor={Research notebook},bookmarksnumbered=true}
\usepackage{bookmark}
\usepackage{tocloft}
\setlength{\cftsecnumwidth}{3em}
\cftsetpnumwidth{2.5em}
\cftsetrmarg{4em}
\usepackage{titlesec}
\titleformat{\section}{\Large\bfseries\raggedright}{\thesection}{0.7em}{}
\usepackage{fancyhdr}
\pagestyle{fancy}\fancyhf{}
\fancyhead[L]{\small\sffamily Open Problems | Research notebook}
\fancyhead[R]{\small 229 | 27 September 2026}
\fancyfoot[C]{\thepage}
\setlength{\parindent}{0pt}
\setlength{\parskip}{5pt plus 1pt}
\setlength{\emergencystretch}{3em}
\setcounter{tocdepth}{1}
\setcounter{secnumdepth}{1}
\setlist[itemize]{leftmargin=3em,itemsep=5pt,topsep=4pt,parsep=0pt}
\allowdisplaybreaks
\newcommand{\N}{\mathbb{N}}
\newcommand{\Z}{\mathbb{Z}}
\newcommand{\Q}{\mathbb{Q}}
\newcommand{\R}{\mathbb{R}}
\newcommand{\C}{\mathbb{C}}
\newcommand{\F}{\mathbb{F}}
\newcommand{\A}{\mathbb{A}}
\newcommand{\PP}{\mathbb P}
\newcommand{\E}{\mathbb{E}}
\newcommand{\eps}{\varepsilon}
\newcommand{\dd}{\,\mathrm d}
\newcommand{\sref}[2]{\hyperlink{source:#1:#2}{[S#2]}}
\newcommand{\eref}[2]{\hyperlink{extra:#1:#2}{[E#2]}}
% Turn the next switch off for a shorter reading copy. The quotations preserve
% every exactTarget field, including qualifications omitted from a short summary.
\newif\ifshowcatalogue
\showcataloguetrue
\newcommand{\cataloguescope}[1]{\ifshowcatalogue
 \paragraph{Exact scope in the supplied catalogue.}
 \begin{quote}\small #1\end{quote}\fi}

% Independently compilable extraction; original section text retained.
\numberwithin{equation}{section}
\begin{document}
\setcounter{section}{0}
% ================================================================
% problemId: 30006801
% Canonical problem ID: 30006801
% exactTarget SHA-256: 64faadfe3be1c0c0d2b566e3c0f25200dfeeb709cdf0ea20970f946f9f2e6fcf
\clearpage
\section*{\texorpdfstring{Berry random-wave conjecture: real negative-curvature BS formulation}{Berry random-wave conjecture: real negative-curvature BS formulation}}\label{30006801}
\subsection{Definitions and mathematical statement}
Let $\mu=\operatorname{vol}_g/\operatorname{vol}_g(M)$ on a closed connected negatively curved $d$-manifold, and choose any real orthonormal eigenbasis with
$\Delta_g\phi_j=\alpha_j^2\phi_j$, $\alpha_j>0$, ordered by eigenvalue. Sample $p$ from $\mu$. The target is weak convergence of the laws of the pointed decorated manifolds
\[
 [M,\alpha_j^2g,p,\phi_j]\ \Longrightarrow\ [\R^d,g_{\mathrm E},0,F]
\]
in the source's local smooth topology on pointed isometry classes. The limiting real Gaussian field has mean zero and covariance
$\E[F(x)F(y)]=\int_{S^{d-1}}\cos\langle x-y,\theta\rangle\dd\sigma(\theta)$,
where $\sigma$ is probability surface measure. Thus it has unit variance and unit spectral radius. Convergence is along the full positive-energy sequence for every real eigenbasis, not just a density-one subsequence or a single moment test.
\par\smallskip\textit{Formulation and concept references:} \sref{30006801}{1}.
\cataloguescope{For every closed connected smooth Riemannian manifold (M,g) of dimension d>=2 with everywhere strictly negative sectional curvature, put mu=dvol\_g/vol\_g(M). For any real orthonormal eigenbasis of L2(M,mu) for the positive Laplace-Beltrami operator, order the eigenvalues nondecreasingly and omit the finite zero-eigenspace; write Delta\_g phi\_j=alpha\_j\textasciicircum{}2 phi\_j with alpha\_j>0. If p is sampled with law mu, does the law of the pointed real-decorated manifold [M,alpha\_j\textasciicircum{}2 g,p,phi\_j] converge weakly, along the full positive-energy sequence, in Abert-Bergeron-Le Masson's local smooth topology on pointed isometry classes of real-decorated d-manifolds, to the law of [R\textasciicircum{}d,g\_Euclidean,0,F]? Here F is the centered real stationary isotropic Gaussian field with covariance E[F(x)F(y)]=integral\_\{S\textasciicircum{}\{d-1\}\} cos(<x-y,theta>) d sigma(theta), where sigma is probability surface measure on the unit sphere (unit variance and unit spectral radius, no extra 2*pi factor).}
\subsection{Short English statement}
When viewed at wavelength scale around a random point, must every high-energy eigenfunction on a negatively curved manifold statistically resemble the universal real Gaussian random wave?
% BEGIN RESEARCH ATTEMPT 229
\subsection{Research attempt}

\paragraph{Current frontier and exact dependency.}
The source's Theorem~4 proves that the proposed Benjamini--Schramm
random-wave convergence implies configuration-space quantum unique
ergodicity: the measures $\phi_j^2\mu$ converge to $\mu$. Its
Sections~2 and~4.2 fix the decorated-space topology and distinguish
convergence of laws from convergence of unbounded moments.
\sref{30006801}{1}. This is an important dependency, not a proof of the
random-wave assertion from quantum ergodicity. A density-one subsequence,
a spectral-window average, and a random linear combination of
eigenfunctions all have different quantifiers from this entry.
The 2024--2026 search in this pass did not produce a verified theorem
covering the unrestricted, full-sequence assertion used here; this is
not an exhaustive audit of every recent manuscript.

\paragraph{Attack route.}
The first route is compactness followed by identification of the limiting
field. A second tries to identify that field from its covariance and the
Helmholtz equation. A third seeks a central-limit mechanism for angular
waves. These routes have distinct missing steps. Compactness alone need
not preserve the full $L^2$ mass; covariance does not characterize a
non-Gaussian law; and classical mixing does not by itself give independent
phases inside one deterministic eigenfunction. The attempt first proves
the compactness calculation and then constructs an exact obstruction to
the second route. The remaining identification problem is formulated
using bounded characteristic functions, avoiding a hidden moment limit.

\paragraph{Attempt.}
\textbf{1. Random-frame tightness at wavelength scale.}
For each $p$, sample a $g$-orthonormal frame $u:\R^d\to T_pM$ with Haar
probability on the orthogonal group, in addition to $p\sim\mu$. Put
\[
 \Psi_j^{p,u}(x)=\phi_j\bigl(\exp_p(u x/\alpha_j)\bigr).
\]
Completeness makes this a smooth function on all of $\R^d$. On each fixed
ball, for all sufficiently large $j$ it is a genuine normal-coordinate
representative of the rescaled decorated manifold. The pulled-back
metric of $\alpha_j^2g$ converges to the Euclidean metric in every
$C^k$ norm, uniformly in $p,u$ on that ball. This follows from the smooth
normal-coordinate metric on the compact frame bundle and
$\alpha_j\to\infty$.

Here is a direct averaged estimate. With $V=\operatorname{vol}_g(M)$,
uniform normal-coordinate Jacobian bounds give, for fixed $R$,
\begin{align}
 \E_{p,u}\int_{B_{R+1}}|\Psi_j^{p,u}(x)|^2\dd x
 &\leq \frac{C_R\alpha_j^d}{V}
       \int_M\int_{B_g(p,(R+1)/\alpha_j)}|\phi_j(y)|^2
                                      \dd\operatorname{vol}_g(y)
                                      \dd\operatorname{vol}_g(p)
       \notag\\
 &\leq C_R'.\label{30006801:average}
\end{align}
The last step is Tonelli followed by the uniform bound on small-ball
volumes and $\int_M|\phi_j|^2\dd\operatorname{vol}_g=V$.
All integrands are nonnegative. The equation in these coordinates is
$\Delta_{g_{j,p,u}}\Psi_j=\Psi_j$, with the positive-Laplacian
convention. Uniform interior elliptic estimates and Sobolev embedding
therefore imply, for every fixed $R,k$,
\begin{equation}\label{30006801:regularity}
 \|\Psi_j^{p,u}\|_{C^{k+1}(B_R)}
       \leq C_{R,k}\|\Psi_j^{p,u}\|_{L^2(B_{R+1})},\qquad
 \E\|\Psi_j\|_{C^{k+1}(B_R)}^2\leq C_{R,k}'.
\end{equation}
One may use intermediate nested balls in the standard elliptic bootstrap;
there is no derivative estimate at the boundary of the larger ball.
This is the same compactness mechanism used in the source's local-limit
framework, especially its Section~8. \sref{30006801}{1}.

Choose a bound for each pair of positive integers $R,k$ so that Markov's
inequality makes the probability of exceeding it less than
$\eta 2^{-R-k}$. Their intersection has probability at least $1-\eta$.
Uniform bounds one derivative higher give relatively compact subsets in
all the lower $C^k(B_R)$ norms, by Arzel\`a--Ascoli and diagonal extraction.
Taking the closure yields a compact subset of $C^\infty_{\rm loc}(\R^d)$.
The finitely many initial $j$ can be included separately, since each is
a continuous image of the compact frame bundle. This proves tightness
of the framed field laws.

Every subsequential law is supported on real solutions of
\begin{equation}\label{30006801:helmholtz}
 -\sum_{a=1}^d\partial_a^2F=F.
\end{equation}
It is invariant under orthogonal changes of coordinates. At the root,
$\E\Psi_j(0)=0$ and $\E\Psi_j(0)^2=1$. The second identity gives uniform
integrability of the \emph{first} power, so the limiting mean is zero;
Fatou's lemma gives only $\E F(0)^2\leq1$. It does not give equality
without uniform integrability of the squares. The source discusses this
possible loss of energy explicitly in its limit-decomposition argument.
\sref{30006801}{1}, Section~4.2. Notably, negative curvature was not used in
\eqref{30006801:average}--\eqref{30006801:regularity}. The difficult geometric input
has not been obtained merely by proving compactness.

\textbf{2. Covariance and the wave equation do not identify the law.}
Let $\Theta$ be uniform on $S^{d-1}$ and $U$ uniform on $[0,2\pi]$,
independent of $\Theta$, and set
\begin{equation}\label{30006801:onewave}
 F_1(x)=\sqrt2\cos(\langle x,\Theta\rangle+U).
\end{equation}
This is a real smooth solution of \eqref{30006801:helmholtz}. Rotation
invariance follows from the law of $\Theta$. Conditional on $\Theta$,
translation of $x$ changes $U$ by a constant modulo $2\pi$, preserving its
uniform law; hence the field is stationary. Direct integration gives
\[
 \E F_1(x)=0,\qquad
 \E F_1(x)F_1(y)=\int_{S^{d-1}}
       \cos\langle x-y,\theta\rangle\dd\sigma(\theta).
\]
Thus it has exactly the covariance and unit variance in the target.
Nevertheless
\begin{equation}\label{30006801:fourth}
 \E F_1(0)^4=4\E\cos^4U=\frac32\ne3.
\end{equation}
More directly, $|F_1(0)|\leq\sqrt2$ almost surely, whereas a standard
normal variable has positive probability outside that interval. A bounded
continuous test supported away from $[-\sqrt2,\sqrt2]$ distinguishes the
laws, without any appeal to convergence of fourth moments.

This is a counterexample to the proposed \emph{identification lemma}
that stationarity, isotropy, unit covariance and the differential equation
force Gaussianity. No negatively curved eigenbasis realizing this law has
been constructed here, so it is not a counterexample to Berry's statement.
Even deriving the covariance for every eigenfunction, rather than for an
average over a spectral window, remains an additional issue.

\textbf{3. The missing central-limit mechanism can be measured exactly.}
Let $F_1^{(1)},\ldots,F_1^{(m)}$ be independent copies of
\eqref{30006801:onewave}, and put $F_m=m^{-1/2}\sum_{a=1}^mF_1^{(a)}$.
It still has the same covariance and solves the same equation. Expanding
the fourth power and using independence, centering and unit variance gives
\[
 \E F_m(0)^4
 =m^{-2}\left(m\frac32+6\binom m2\right)
 =3-\frac{3}{2m}.
\]
The finite-dimensional multivariate central limit theorem makes these
particular random sums converge to the Gaussian field's finite-dimensional
laws as $m\to\infty$. This is a model calculation, not a construction of
the prescribed $\phi_j$. For one deterministic eigenfunction, neither an
independent-phase representation nor a growing number of effectively
independent waves has been proved in this notebook. Invoking mixing of
the geodesic flow without supplying that bridge would assume the main
probabilistic step rather than establish it.

\textbf{4. A precise bounded-test completion criterion.}
Write $K(z)=\int_{S^{d-1}}\cos\langle z,\theta\rangle\dd\sigma(\theta)$.
In conjunction with the proved tightness, it would suffice to establish,
for every finite collection $x_1,\ldots,x_s\in\R^d$ and
$t_1,\ldots,t_s\in\R$,
\begin{equation}\label{30006801:characteristic}
 \lim_{j\to\infty}\E_{p,u}
    \exp\left(i\sum_{a=1}^s t_a\Psi_j^{p,u}(x_a)\right)
 =\exp\left(-\frac12\sum_{a,b=1}^s t_at_bK(x_a-x_b)\right).
\end{equation}
Indeed every subsequential smooth law would then have the same
finite-dimensional characteristic functions as the required Gaussian law.
Evaluations on a countable dense set determine a continuous function and
generate the Borel sigma algebra on this smooth-function space; thus the
law is unique. Tightness and uniqueness of all subsequential limits imply
convergence of the whole sequence. Passing from frames to pointed
isometry classes is continuous in the source's local smooth topology,
so this stronger framed assertion gives the exact unframed conclusion.
The integrands in \eqref{30006801:characteristic} are bounded by one, so no
unproved uniform integrability is hidden in this reduction.

\paragraph{Stress test.}
The single-wave construction breaks the covariance route even after all
its proposed regularity and symmetry hypotheses are granted. It does not
break the conjecture, whose eigenfunction origin imposes information not
contained in that construction. Conversely, rare large peaks and divergent
fourth moments alone need not break weak convergence to a Gaussian field:
unbounded moment tests are not automatically admissible in the stated
topology. The source makes that distinction in discussing arithmetic
examples. \sref{30006801}{1}, Section~4.2.

Every estimate in the tightness proof has constants depending on the fixed
manifold, radius and derivative order, which is permitted. It does not
assert uniformity over all negatively curved manifolds. More importantly,
\eqref{30006801:characteristic} is required along the \emph{full} sequence and
for every real eigenbasis, including choices inside eigenspaces. Replacing
it by a density-one, random-basis or energy-averaged statement changes the
target. Its configuration-space QUE consequence is a warning about its
strength, not a reason to stop attempting it.

\paragraph{Outcome.}
\textbf{Outcome: Exploratory approach with unresolved gap.}

The notebook gives a direct random-frame tightness argument, an exact
non-Gaussian stationary field with the prescribed covariance, and a
bounded-characteristic-function criterion that would complete the proof.
It also computes how the Gaussian fourth moment emerges in an independent
superposition model. These are deductions and diagnostics, not claims of
a new negative-curvature theorem or priority for the compactness method.

The unresolved step is identification of the deterministic eigenfunction
limits, specifically a proof of \eqref{30006801:characteristic} or an equally
strong substitute using negative curvature. No such estimate and no
counterexample eigenbasis is established here.
% END RESEARCH ATTEMPT 229
\subsection{Sources}
\begin{itemize}\raggedright
\item[\textnormal{[S1]}]\hypertarget{source:30006801:1}{} Eigenfunctions and random waves in the Benjamini-Schramm limit; arXiv1810.05601v2.
\url{https://arxiv.org/pdf/1810.05601v2}.
{\small\textit{Catalogue formulation source.}}
% BEGIN ADDED REFERENCES 229
% No additional external reference is needed: the consulted formulation
% source [S1], especially Sections 2, 4.2 and 8 and Theorem 4, is cited locally.
% END ADDED REFERENCES 229
\end{itemize}

\end{document}
